\documentclass[../../../include/open-logic-section]{subfiles}

\begin{document}

\olfileid{sth}{replacement}{ref}
\olsection{Panggantian lan Refleksi}

Upaya pungkasan kita kanggo menehi dhasar
marang Panggantian liwat \stagesinex diwiwiti
kanthi asil kang jero lan endah:\footnote{Elinga:
kabeh formula bisa ngemot parameter (kajaba yen
diandharake kanthi cetha kosokbaline).}
%Ing kene aku bakal nganggo garis ndhuwur, kayata $x_1, \ldots, x_n$, kanggo
%nyingkat ``$x_1, \ldots, x_n$'':
\begin{thm}[Skema Refleksi]\ollabel{reflectionschema}
Kanggo saben formula $\phi$:
\[
\forall \alpha \exists \beta > \alpha (\forall x_1 \ldots, x_n \in
V_\beta)(\phi(x_1, \ldots, x_n) \liff \phi^{V_\beta}(x_1, \ldots, x_n))
\]
\end{thm}
\noindent 
Kaya ing \olref[sth][replacement][strength]{formularelativization},
$\phi^{V_\beta}$ dipikolehi kanthi matesi
saben kuantor ing $\phi$ marang himpunan~$V_\beta$.
Dadi, miturut intuisi, Refleksi kandha:
yen $\phi$ bener ing sakabehing hierarki,
$\phi$ uga bener ing \emph{perangan wiwitan}
hierarki kang indeksé bisa ngluwihi ordinal
apa wae.

\citet{Montague1961} lan \citet{Levy1960}
nuduhake yen (rumusan kang trep saka)
Panggantian lan Refleksi iku ekivalen relatif
marang~$\Z$, mula nambah salah siji bakal
ngasilake~$\ZF$. (Kita bakal mbuktekake iki
ing \olref[sth][replacement][refproofs]{sec}.)
Amarga ekivalensi iki, kita bisa ngarepake
yen Refleksi lan Panggantian bisa diwenehi
dhasar liwat \stagesinex{} mangkene: miturut
\stagesinex, hierarki mesthine dhuwur banget,
nganti pratelan apa wae babagan hierarki ora
cukup kanggo matesi dhuwure. Iki bisa
dimangerteni mangkene: yen sawijining
ukara~$\phi$ bener ing sakabehing hierarki,
ing sadhuwure wates apa wae ana perangan
wiwitan hierarki kang uga ndadekake ukara iku
bener.
Iki persis Refleksi.

Sepisan maneh, iki katon minangka upaya kang
njanjeni kanggo menehi dhasar intrinsik marang
Panggantian. Nanging rembugane akeh banget
kanggo diwenehake ing kene. Pamaca kudu
mutusake dhewe apa upaya iki kasil.\footnote{Kanggo
nerusake, pamaca bisa maca
\citet[95--100]{Incurvati2020}.}

\end{document}
